% !TeX root=../main.tex

\chapter{ابرپارامترها و جزئیات پیاده‌سازی}
\label{app:hyperparams}

این پیوست، جدول کامل ابرپارامترهای استفاده‌شده در تمام آزمایش‌های این پایان‌نامه را، تکمیل‌کنندهٔ بخش~\ref{sec:impl} از فصل سوم، فهرست می‌کند.

\begin{table}[htb]
\centering
\caption{ابرپارامترهای مدل، لورا و بهینه‌سازی.}
\begin{tabular}{ll}
\toprule
پارامتر & مقدار \\
\midrule
مدل پایه & \lr{Pythia-70m-deduped} \\
تعداد لایه‌ها & ۶ \\
ابعاد جریان باقیمانده & ۵۱۲ \\
رتبهٔ لورا ($r$) & ۸ \\
ضریب مقیاس لورا ($\alpha$) & ۱۶ \\
دراپ‌اوت لورا & ۰۵٫۰ \\
ماژول‌های هدف لورا & \lr{query\_key\_value, dense} \\
بهینه‌ساز & \lr{AdamW} \\
نرخ یادگیری & \lr{$10^{-4}$} \\
تعداد گام‌های آموزش & ۵۰۰ \\
بازهٔ ثبت چک‌پوینت & هر ۵۰ گام \\
کران بریدن نرم گرادیان & ۰٫۱ \\
\bottomrule
\end{tabular}
\end{table}

\begin{table}[htb]
\centering
\caption{ابرپارامترهای دادگان.}
\begin{tabular}{ll}
\toprule
پارامتر & مقدار \\
\midrule
اندازهٔ دستهٔ \lr{WikiText} در هر گام & ۸ دنباله \\
طول هر دنبالهٔ \lr{WikiText} & ۲۵۶ نشانه \\
اندازهٔ دستهٔ \lr{PAWS} در هر گام & ۸ جفت هم‌معنی + ۸ جفت ناهم‌معنی \\
جفت‌های آموزش \lr{PAWS} (هر شرط) & ۴۰۰۰ \\
جفت‌های اعتبارسنجی \lr{PAWS} (هر شرط) & ۳۰۰ \\
جفت‌های آزمون \lr{PAWS} (هر شرط) & ۱۰۰۰ \\
\bottomrule
\end{tabular}
\end{table}

\begin{table}[htb]
\centering
\caption{ابرپارامترهای تابع هزینه و ارزیابی.}
\begin{tabular}{lp{7.5cm}}
\toprule
پارامتر & مقدار \\
\midrule
وزن منظم‌ساز بازسازی خودرمزگذار & ۰٫۱ \\
کران بالای خطای بازسازی نرمال‌شده (رابطهٔ~\ref{eq:recon}) & ۱۰ \\
مقادیر آزموده‌شدهٔ $\lambda$ (شرایط \lr{C2}--\lr{C4}) & $\{0.01,\ 0.03,\ 0.1,\ 0.3,\ 1.0\}$ \\
بازهٔ مقیاس آموختنی \lr{BCE} ($\gamma$) & $[1,\ 50]$ \\
آستانهٔ هشدار اعتبار خودرمزگذار & $90\%$ \\
لایه(های) هدف آزموده‌شده & $\{0\},\{1\},\{2\},\{3\},\{4\},\{5\},$ $\{2,3\},\{3,4\},\{3,4,5\},\{4,5\}$ \\
\bottomrule
\end{tabular}
\end{table}

جدول~\ref{tab:app-grid} کل شبکهٔ آزمایشی را--با تفکیک شرط و تعداد اجراهای هر شرط--خلاصه می‌کند.

\begin{table}[htb]
\centering
\caption{ساختار کامل شبکهٔ آزمایشی. هر سطر، تعداد اجراهای مربوط به آن شرط را در \emph{هر} یک از ده حالت انتخاب لایه نشان می‌دهد؛ عدد کل، حاصل‌ضرب این مقدار در ۱۰ است.}
\label{tab:app-grid}
\begin{tabular}{lccc}
\toprule
شرط & تعداد مقادیر $\lambda$ & اجرا به ازای هر حالت لایه & کل اجراها (در ۱۰ حالت) \\
\midrule
\lr{C1\_lm\_only} & --- (بدون $\lambda$) & ۱ & ۱۰ \\
\lr{C2\_raw\_invariance} & ۵ & ۵ & ۵۰ \\
\lr{C3\_sae\_magnitude} & ۵ & ۵ & ۵۰ \\
\lr{C4\_sae\_magnitude\_support} & ۵ & ۵ & ۵۰ \\
\midrule
جمع اجراهای آموزشی & & & ۱۶۰ \\
خط پایه (بدون ریزتنظیم) & & & ۱ (مشترک) \\
\bottomrule
\end{tabular}
\end{table}

\paragraph{یادداشتی دربارهٔ پیکربندی توصیه‌شده.}
همان‌طور که در فصل چهارم استدلال شد، پیکربندی توصیه‌شدهٔ این پایان‌نامه--شرط
\lr{C4}، لایهٔ هدف منفرد پنجم، و
$\lambda=0.3$--لزوماً همان پیکربندی‌ای نیست که بالاترین بهبود خام
\lr{AUROC}
را در کل شبکهٔ آزمایشی به دست می‌دهد؛ این انتخاب بر مبنای مصالحه‌ای آگاهانه میان بهبود خام، سرگشتگی زبانی، و اعتبار بازسازی خودرمزگذار صورت گرفته است، نه صرفاً بیشینه‌سازی یک سنجهٔ منفرد.
